\section*{الفصل \ref{ch:g1:magnets} --- المغانط}

\begin{solution}{exo:g1:magnets:1}
يمسك بها: المسمار، ومشبك الورق، والملعقة الفولاذية. ويهمل:
الفلّينة، والبطّة البلاستيكية، والورقة.
\end{solution}

\begin{solution}{exo:g1:magnets:2}
الحديد (والفولاذ، وأكثره حديد). ويُهمَل مع ذلك: قطعة النقود
ورقاقة الألمنيوم --- معدنان، لكنهما ليسا حديدًا.
\end{solution}

\begin{solution}{exo:g1:magnets:3}
الورقة (الرسم) بينهما. وجذب \omterm{def:g1:magnets:magnet}{المغناطيس} ينفذ من الورقة نفاذًا إلى
الباب الحديدي --- فهو يعمل من خلال الأشياء.
\end{solution}

\begin{solution}{exo:g1:magnets:4}
أَمِلَّ \omterm{def:g1:magnets:magnet}{المغناطيس} على \emph{الوجه الخارجي} للكأس، بجوار المشبك،
وامشِ به صاعدًا على جدارها على مهل: يتبعه المشبك من الداخل حتى
الحافة، مجذوبًا عبر الزجاج والماء.
\end{solution}

\begin{solution}{exo:g1:magnets:5}
بعيدًا عن \omterm{def:g1:magnets:magnet}{المغناطيس} يكون الجذب أضعف من أن يحرّك المشبك؛ ولا
يشتدّ إلا عن قرب. فمدى \omterm{def:g1:magnets:magnet}{المغناطيس} يخبو سريعًا كلما زادت المسافة.
\end{solution}

\begin{solution}{exo:g1:magnets:6}
إدارة أحد المغناطيسين، حتى يتقدّم طرفه الآخر. فعندئذ يتدافعان
بدل أن يلتصقا.
\end{solution}

\begin{solution}{exo:g1:magnets:7}
عند الطرفين --- إذ تهمل المشابك الوسط في الغالب. \omterm{def:g1:magnets:magnet}{فالمغناطيس} أقوى
ما يكون عند طرفيه.
\end{solution}

\begin{solution}{exo:g1:magnets:8}
وجه الشبه: كلاهما يستطيع أن يسحب شيئًا نحوك. ووجه الاختلاف: يدك
لا بدّ أن تلمس اللعبة، أما \omterm{def:g1:magnets:magnet}{المغناطيس} فيسحب المشبك عبر الهواء
الخالي، ومن خلال الخشب والماء، ولا يلمس شيئًا.
\end{solution}

\begin{solution}{exo:g1:magnets:9}
قرّب \omterm{def:g1:magnets:magnet}{المغناطيس} من كل جسم على حدة. يقفز المسمار إليه؛ ولا تتحرّك
قطعة النقود ولا الرقاقة، مع أن كليهما معدن. فجملة توم خطأ كما
قالها؛ وتصحيحها: ``تجذب المغانط ما صُنع من حديد (أو فولاذ)، لا
كل معدن.''
\end{solution}
